\section{Inductive Logic Programming}
\label{section-popper}
\label{sect-ilp}

Inductive Logic Programming lies at the intersection of machine learning and logic programming. In brief, it aims to learn logical
rules from positive and negative examples, as well as from background knowledge. All of these elements, rules, examples, and background
knowledge are expressed within the framework of logic programming. To provide the reader with the minimum background necessary to understand
our work, we briefly describe what logic programming is, what inductive logic programming is, and, in particular, the principles on
which the Popper language has been designed.

\subsection{Logic Programming}
Logic programming is a programming paradigm based on formal logic. The main idea is that, instead of providing the computer with instructions on how to solve a problem, as in procedural languages like C or Java, one specifies by logical rules the problem to be solved and let execution models execute them so as to actually solve the problem. This idea assumes the existence of two semantics, called declarative semantics and operational semantics, which provide meanings to programs either according to a (declarative) logical viewpoint or to an operational (execution-based) viewpoint. It is the equivalence between the two semantics that allow the user to specify the problem to be solved and let the computer executes the necessary computational steps. To be complete a third component needs to be addressed: the syntax according to which well-formed formulae are written. These three components are the subjects of the following
three subsections.

\subsubsection{Syntax.}
Any first-order logic theory is built upon three classes of symbols: a set of variables, a set of function names together with their arities
and a set of predicate names together with their arities. In the following, we shall use the Prolog syntax:

\begin{itemize}

\item variables are represented as strings of characters starting with an uppercase letter
\item function names and predicate names are represented as strings characters starting with a lower case letter.
  
\end{itemize}

\noindent
The arity of a function name or predicate name is the number of arguments it takes. As a result, $X$ is typically used subsequently to
represent a variable and $f/3$, $p/2$ are respectively examples of function and predicate names.

The three classes of symbols are assumed to be disjoint two by two. Constants are a special case of function names of arity
$0$. Examples are $\mathit{alice}$ and $\mathit{bob}$.

Terms and atoms are built from these classes as follows.

\begin{definition}
Terms are inductively defined as follows:

\begin{itemize}

\item a variable is a term
\item a constant is a term
\item if $f/n$ is a function name of arity $n$, and if $t_1$, \ldots, $t_n$ are terms then
  $f(t_1,\cdots,t_n)$ is a term
  
\end{itemize}
\end{definition}

\begin{definition}
An atom is a construct of the form $p(t_1,\cdots,t_n)$ where $p/n$ is
a predicate name of arity $n$ and $t_1$, \ldots, $t_n$ are terms. A
literal is an atom or the negation of an atom.
\end{definition}

We are now in a position to define the logical rules to be used. Those are named clauses.

\begin{definition}
A clause is a universally quantified disjunction of the form:
\[ \forall X_1 \cdots \forall X_n \; L_1 \vee \cdots \vee L_m \]
where $L_1$, \ldots, $L_m$ are literals and $X_1$, \ldots, $X_n$ are all the variables appearing in them. 
\end{definition}

Very often the universal quantification is taken as implicit and is not
written.  Special forms are worth mentionning.

\begin{definition}
A Horn clause is a clause containing at most one positive literal. A
definite clause is such a clause that contain one positive
literal. By using classical reasoning, it is easy to verify that
it can be rewritten as the following implication:
\[ h \leftarrow b_1 \wedge \cdots \wedge b_n \]
with $h$, $b_1$, \ldots $b_n$ (positive) atoms.
A fact is a definite clause that contains no atoms $b_i$. It is simply rewritten as \( h \).
A goal is a clause that contains no positive literals. By similar reasoning, it can be rewritten as
\[ \leftarrow b_1 \wedge \cdots \wedge b_n \]
\end{definition}

The intuition behind definite clauses is that if the body
$b_1 \wedge \cdots \wedge b_n$ is true, then the head $h$ is true. Facts are then atoms
that are always true (without any additional assumption). Finally,
goals are atoms that need to be established.

It will be useful later to give values to variables. This is achieved by the notion of substitution.

\begin{definition}
A substitution is a finite set of the form
\( \{ X_1/t_1, \cdots, X_n/t_n \} \)
where $X_1$, \ldots, $X_n$ are distinct variables,
$t_1$, \ldots, $t_n$ are terms, and where $t_i$ is distinct from $X_i$,
for any $i$.
\end{definition}

Substitutions are subsequently denoted by greek letters, such as $\theta$ and $\sigma$.
The intuition behind substitutions is that each variable $X_i$ is
given $t_i$ as a value. This leads to the application of substitutions
to terms, atoms or even clauses.

\begin{sloppypar}
\begin{definition}
Let $tac$ be either a term, an atom or a clause and let
$\theta= \{ X_1/t_1, \cdots, X_n/t_n \}$
be a substitution. Then the application of $\theta$ to $tac$,
subsequently denoted $tac \, \theta$, is the term, atom or clause,
obtained by simultaneously replacing each occurrence of $X_i$ by
$t_i$.
\end{definition}
\end{sloppypar}
  
This notion finds a direct application in unification.

\begin{definition}
Two terms (or atoms) $ta_1$ and $ta_2$ unify if there is a
substitution $\theta$ such that \( ta_1\theta = ta_2\theta \).  In
that case, $\theta$ is said to be a unifier of $ta_1$ and $ta_2$.
\end{definition}

Another application consists in applying subsitutions
successively. This leads to composing substitutions.

\begin{definition}
Let $\theta = \{ X_1/t_1, \cdots, X_m/t_m \}$ and
$\sigma = \{ Y_1/u_1, \cdots, Y_n/u_n \}$
be two substitutions. The composition of $\theta$ by $\sigma$, denoted
$\theta\sigma$, is the substitution obtained from the set
$\{ X_1/(t_1\sigma), \cdots, X_m/(t_m\sigma), Y_1/u_1, \cdots, Y_n/u_n \}$
by removing any binding $X_i/(t_i\sigma)$ for which $X_i = t_i\sigma$
and any binding $Y_j/u_j$ for which $Y_j \in \{ X_1, \cdots, X_m \}$.
\end{definition}

\begin{definition}
The substitution $\theta$ is said to be more general than the
substitution $\sigma$ if there exists a substitution $\gamma$ such
that $\sigma = \theta\gamma$.
\end{definition}

\begin{definition}
Let $t$ and $u$ be two terms or atoms. The substitution $\theta$ is
said to be a more general unifier (mgu for short) of $t$ and $u$ if it
is a unifier of $t$ and $u$ and if it is more general than any other
unifier.
\end{definition}

\subsubsection{Declarative Semantics.}
The declarative semantics of logic programs defines their meaning
purely in logical terms without any reference to any execution
model. It is based on the notions of interpretation and model. A full
treatment of this semantics is out of the scope of this paper. We
refer the interested reader to \cite{Lloyd} for a complete
overview. However the key points may be described on Horn clauses
through the notion of Herbrand universe and Herbrand base.

\begin{definition}
A term, atom or clause is said to be ground if it contains no
variable. The Herbrand universe, denoted $U_H$, is defined as the set
of all ground terms. The Herbrand base, denoted $B_H$, is defined as
the set of ground atoms.
\end{definition}

\begin{definition}
Let $tac$ and $tac'$ be two terms (resp. atoms or clauses). We say
that $tac'$ is a ground instance of $tac$ iff $tac'$ is ground and
there is a substitution $\theta$ such that $tac' = tac\theta$. In the
following, we shall denote by $\mathit{ground(tac)}$ the set of the
ground instances of $tac$.
\end{definition}


\begin{definition}
An Herbrand interpretation is a subset of the Herbrand base $B_H$.
\end{definition}

The intuition behind an Herbrand interpretation is that it defines as
true the atoms it contains and as false the other ground atoms. Then,
by using the classical truth tables, it is easy to determine whether a
given ground clause is true. This allows us to define the notion of model.

\begin{definition}
A Herbrand interpretation $I$ is a model of a Horn clause $c$ iff the
truth value of any ground instance of $c$ determined with respect to
$I$ is true. It is a model of a set of Horn clauses iff it is a model
for each of the clauses of the set.
\end{definition}

\begin{definition}
A goal $G$ is a logical consequence of a set of Horn clauses $S$ iff
every Herbrand model of $S$ is also a Herbrand model of $G$. This is
subsequently denoted as $S \models G$.
\end{definition}

A major result of model theory in the case of Horn clauses is that the
set $M_P$ of logical consequences of a set of Horn clauses $P$ can be determined
as the intersection of all its Herbrand models and may be further
characterized as the least fixed point of the immediate consequence
operator $T_P$ defined as follows.

\begin{definition}
Let $P$ be a set of Horn clauses. The immediate consequence operator
$T_P: {\cal P}(B_H) \rightarrow {\cal P}(B_H)$ is defined as follows: for
any Herbrand interpretation $I \subseteq B_H$,
\[ T_P(I) = \{ h : \begin{array}[t]{l}
                  (h \leftarrow b_1, \cdots, b_n) \mbox{ is a ground instance of a clause of } P \\
                  b_1, \cdots, b_n \in I \}
                 \end{array}
\]
\end{definition}

Obviously, by construction, one has $T_P(M_P) = M_P$. It turns out
that this result can be generalized for Horn clauses extended with
negative literals in their body. We refer the reader to \cite{AptBol} for
more explanations. For this introduction, it sufficient to say that,
under some conditions, the stable semantics developped by M. Gelfond
and V. Lifschitz has allowed to identify stable models in the sense
that they are fixed points of (suitable extensions) of the $T_P$
operator. This has lead to develop so called answer set programming,
which is a framework in which one computes stable models for logic
programs, potentially extended to include constraints and optimization.


\subsubsection{Operational Semantics.}
\begin{sloppypar}
From an operational perspective, the computation proceeds
as a succession of steps that transform a goal (or conjunction of atoms)
into another one. Given a goal, say
\( \leftarrow A_{1}, \cdots, A_{n} \),
one selects non-deterministically an atom, say $A_i$, and then a
definite clause\footnote{%
  When need be, the variables of the clause are previously
  renamed so that none of them appears earlier in the computation.
}
\( H \leftarrow B_{1},  \cdots , B_{k} \)
whose head $H$ unifies with $A_i$.
Let $\theta$ be the most general unifier of $A_{i}$ and $H$.  The
conjunction
\( \leftarrow A_{1}, \cdots, A_{n} \)
is then transformed into the new conjunction 
 \( \leftarrow ( A_{1}, \cdots, A_{i-1}, B_{1},
 \cdots, B_{k}, A_{i+1}, \cdots, A_{n} ) \theta.  \)
 Such a transformation step is called a \textit{reduction step} and any
 sequence of reduction steps is called a \textit{SLD-derivation}.  The
 reduction process ends when the current conjunction is empty or when
 no reduction step can take place.
 \end{sloppypar}
   
SLD-derivation can be described more formally by means of the derivation relation
\( \jmjderiv{P}{\mfm{G}}{\theta} \).
 Its intended meaning is that there is a (successful) {\em SLD-derivation of\/} \mf{G} with respect to \mf{P}
 yielding the computed answer substitution $\theta$. This derivation relation is
 itself defined by means of rules of the form
 \[\frac{\mfm{Assumptions}}{\mfm{Conclusion}} \hspace*{4ex}
    \mbox{\textit{if}}
 \hspace*{2ex} \mfm{Conditions},  \]
 %
asserting the \mf{Conclusion} whenever the \mf{Assumptions} and
\mf{Conditions} hold. Note that \mf{Assumptions} and \mf{Conditions}
 may be absent from some rules.  Precisely, it is defined, by case
 analysis on the form of \mf{G} as the smallest relation satisfying 
 rules (E) and (A) of Figure~\ref{jmj-figderiv}.

 \begin{figure}[t]
%
\begin{center}
% \begin{tabular}{|c|}
\begin{tabular}{c}  
% \hline
%
\begin{tabular}{lll}
\\
  (E) & \hspace*{2ex} &
   $\begin{array}{c}
      \\ 
         \hline 
         \jmjderiv{ \mfm{P} }{ \jmjemptyconj }{ \epsilon }
   \end{array}$
\\ \mbox{ } 
\\  \mbox{ } 
\\
  (A) & \hspace*{2ex} &
   $\begin{array}{c}
         \jmjderiv{ 
                     \mfm{P} 
                  }{
                     ( \mfm{A}_{1}, \cdots, \mfm{A}_{i-1},
                       \mfm{B}_{1}, \cdots, \mfm{B}_{n},
                       \mfm{A}_{i+1}, \cdots, \mfm{A}_{m} ) \theta
                       \; \; 
                  }{
                     \sigma
                   }
      \\ \hline 
         \jmjderiv{ 
                     \mfm{P} 
                  }{ 
                     \mfm{A}_{1}, \cdots, \mfm{A}_{m}
                     \; \;
                  }{
                     (\theta\sigma)_{\mid \mfm{ 
                               vars( \bsetm A_{1},\cdots,A_{m} \esetm ) }}
                   }
   \end{array}$
\\ \mbox{ }
\\ & & 
\hspace*{6ex} {\em if } 
\( \begin{system}{l}
      ( \mfm{H} \leftarrow \mfm{B}_{1}, \cdots, \mfm{B}_{n} ) 
      \mbox{ is a fresh renaming of a clause in } \mfm{P} 
      \\
      \mfm{A}_{i} \mbox{ and } \mfm{H} \mbox{ unify with mgu } \theta
\end{system} \)
\\  \mbox{ }
\\  
\end{tabular}
%
\\
%\hline
\end{tabular}
\end{center}
%
\caption{The derivation relation
\label{jmj-figderiv}}
\end{figure}

Rule (E) states that the empty conjunction is derivable from any
program with the empty substitution $\epsilon$ as computed answer
substitution.  Rule (A) re-expresses the reduction step described
previously.

The equivalence between the declarative and operational semantics is
an attractive feature of logic programming.  It allows the programmer
to concentrate on the declarative aspects of the programming task and
let the execution model worry about operational details. This result
is stated by the following theorem, for which a proof can be found in
\cite{Lloyd}.


\begin{thm}
Let $P$ be a set of definite clauses and $G$ be a goal.

\begin{thmprop}

\item
{\em Soundness.\/}
For any substitution $\theta$, if 
\( \jmjderiv{ \mfm{P} }{ \mfm{G} }{ \theta } \)
holds, then so does
\( \jmjlogcons{ \mfm{P} }{ \mfm{G} }{ \theta } \).

\item
{\em Completeness.\/}
For any substitution $\theta$ such that
\( \jmjlogcons{ \mfm{P} }{ \mfm{G} }{ \theta } \)
holds, there are two substitutions $\gamma$ and $\mu$ such that 
\( \jmjderiv { \mfm{P} }{ \mfm{G} }{ \mu } \)
and
\( \mfm{G}\theta = \mfm{G}\mu\gamma \).

\end{thmprop}
\end{thm}


\subsection{Inductive Logic Programming}

Inductive logic programming (ILP for short)~\cite{muggleton1995inverse} aims to learn a hypothesis $H$
describing a predicate as a logic program, using positive examples $E^+$ of this predicate, negative
examples $E^-$ of this predicate, and a background knowledge $B_k$, and this in a way such that
$H \cup B_k$ entails all positive examples and none of the negative ones:
%
\begin{eqnarray} 
B_k \cup H \models e^+ \mbox{ for any } e^+ \in E^+ & & \label{posentail} \\
B_k \cup H  \nvDash e^- \mbox{ for any } e^- \in E^- & \label{negentail}
\end{eqnarray}

\noindent
ILP is typically applied in a centralized setting 
in which all examples are available on a single machine.

To further characterize ILP, several concepts need to be introduced. 
The background knowledge $B_k$ is used to provide a context for
learning. It contains a set $B$ of facts and rules that describe
relationships already understood about the domain. It also includes a
so-called declaration biais $D$, which aims at restricting the
hypotheses that can be constructed. Examples include predicate
declarations, which specify which predicate names can appear in the
head of clauses, in particular the form of the arguments, as well as body
declarations, which do the same for predicates appearing in the body
of clauses. It is assumed that no predicate symbol in the body of a
clause in $B_k$ appears in a head declaration of D. In other words, we
assume that $B_K$ does not depend on any hypothesis.

\begin{sloppypar}
It is worth noting that predicate declarations can be expressed as constraints at the
level of a meta-language. This is one of the core properties of the
language Popper~\cite{popper} we shall use. Assuming such a meta-language ${\cal L}$, 
a hypothesis constraint is a constraint expressed in ${\cal L}$. For instance, 
the atom \texttt{head\_literal(Clause,Pred,Arity,Vars)} is used in Popper to denote that the clause
\texttt{Clause} has a head literal with the predicate symbol \texttt{Pred}, is of arity \texttt{Arity}, and has the
arguments \texttt{Vars}. As a further use
\end{sloppypar}

\begin{lstlisting}
  :- head_literal(_,p,2,_).
\end{lstlisting}

\noindent
states that a predicate symbol \texttt{p} of arity 2 cannot appear in the head of
any clause in a hypothesis.

The input of the inductive learning problem can then formulated as composed of the following ingredients.

\begin{definition}
\label{ilp-problem}
The ILP problem input is a tuple $(B, D, C, E^+, E^-)$ where

\begin{itemize}

\item $B$ is a Horn program denoting background knowledge
\item $D$ is a declaration biais
\item $C$ is a set of hypothesis constraints
\item $E^+$ is a set of ground atoms denoting positive examples of the predicate to be learnt
\item $E^-$ is a set of ground atoms denoting negative examples of the predicate to be learnt.

\end{itemize}
\end{definition}


The entailment relations~\ref{posentail} and~\ref{negentail} characterize complete and consistent hypotheses.

\begin{definition}
Let $(B, D, C, E^+, E^-)$ be an input tuple and $H$ be a hypothesis. Then $H$ is:

\begin{itemize}
    \item complete when \( B \cup H \models e^+\) for any  $e^+ \in E^+$
    \item consistent when \( B_k \cup H  \nvDash e^- \) for any $e^- \in E^-$ 
\end{itemize}

\noindent
A hypothesis is a solution when it is both complete and consistent.
\end{definition}

In practice there can be many solutions to an inductive learning
problem input. One is thus naturally interested by optimal solutions,
which are typically defined by minimizing the number of literals used.

\begin{definition}
The size of a hypothesis $H$, denoted $size(H)$ is the total number of literals
employed in $H$. Given an inductive learning problem input, an optimal solution is defined as
a solution $H$ such that any other solution $H'$ verifies $size(H) \leq size(H')$.
\end{definition}

It is not always possible to find a complete and consistent
hypothesis. In that case, one tries to get the best hypothesis in the
sense that it covers the most positive examples and the less negative examples. This leads to the notions of confusion matrix and score.

\begin{definition}
\label{confusion-matrix}
Let $(B, D, C, E^+, E^-)$ be an input tuple and $H$ be a
hypothesis. Then the confusion matrix for $H$, denoted
$conf\_matrix(H)$, is the tuple $(TP,FN,TN,FP)$ whose components are defined as follows:

\begin{itemize}

\item $TP$ is the number of positive examples correctly entailed by $B \cup H$ (true positives) 
\item $FN$ is the number of positive examples not entailed by $B \cup H$ (false negatives) 
\item $TN$ is the number of negative examples not entailed by $B \cup H$ (true negatives)
\item $FP$ is the number of negative examples incorrectly entailed by $B \cup H$ (false positives)

\end{itemize}

\noindent
The score of $H$, denoted $score(H)$, is defined as $TP + TN$.
\end{definition}

In the past thirty years, many inductive logic languages have been introduced. The reader is referred to article~\cite{Cropper-30} for a detailed description. Pioneer languages, such as Progol and Metagol, have explored subsumption lattices, amounting to trees of rules. In contrast, Popper has taken a meta-language approach which allows to generate hypotheses by using Answer Set Programming (ASP) and which allows to learn from failures. It is generally considered as very appealing since it allows to find globally optimal hypotheses whereas Progol and Metagol may produce only locally optimal hypotheses. We shall turn in this paper to Popper. Its main algorithm is described in Algorithm~\ref{alg:popper} as a Python function inspired from the one presented in~\cite{cropper2021learning}. The function takes as arguments the sets $E^+$, $E^-$, $D$, $C$ described in Definition~\ref{ilp-problem}. Three other arguments are added to limit the search space: the maximum number of variables, literals and clauses in a hypothesis. The Python function returns a hypothesis: either a solution, if there is one, or an hypothesis that maximizes the score function otherwise. The algorithm mainly consists of two loops. The outer one is responsible for searching hypotheses which increasingly add literals. For a fixed number of literals, the inner loop embodies the Popper generate-test-constraint loop. Given a current set \texttt{CC} of constraints, the call to the \texttt{generate} function generates a hypothesis. This is obtained through Answer Set Programming acting over the constraints \texttt{CC} of a meta-language. The produced solution, if there is one, is a hypothesis meeting the constraints of \texttt{CC}. In case there is no solution then the variable \texttt{c\_solvable} is set to false which leads to finish the inner loop. Otherwise the solution is tested against the positive and negative examples by additionally using the background knowledge of \texttt{B}. This is operated by the call to the function \texttt{test}. This results in a confusion matrix as specified in Definition~\ref{confusion-matrix}. Two outputs are computed therefrom: an \texttt{outcome} and a \texttt{score}. The outcome summarizes how many positive and negative examples are covered. For the positive examples, either all of them are covered (which is denoted as ``all''), or  some are covered (which is denoted by ``some'') or none are covered (which is denoted by ``none''). For the negative examples, either some are covered (which is denoted by ``some'', this including the case where all are covered) or none are covered (which is denoted by ``none''). Obviously a solution is found in case (all,none) is produced. In the other cases, the article~\cite{cropper2021learning} explains which constraints are to be added to the \texttt{CC} constraints. Note that this is operated so as to exclude answer sets in the ASP resolution, which makes the algorithm converge. Technically in the algorithm this is achieved by the call to the function \texttt{learn\_constraints}. The important point for our research is that the added constraints can be determined in view of the outcome and do not need the knowledge of positive nor negative examples. It is also worth noting that, having computed its score, the new hypothesis is retained as the best one if its score is better than the one memorized in variable \texttt{best\_hyp}. 



\begin{algorithm}[t]
  \caption{Popper (inspired by \cite{cropper2021learning})}
  \label{alg:popper}
  \small
\begin{lstlisting}
def popper(E+, E-, B, D, C, max_vars, max_literals, max_clauses):

  CC = C
  best_score = 0
  best_hyp = { }
  end_popper = false
  num_literals = 0

  while not end_popper and num_literals <= max_literals:
    c_solvable = true
    num_literals += 1
    while not end_popper and c_solvable:
       hyp = generate(D, CC, max_vars, num_literals, max_clauses)
       if not hyp: 
         c_solvable = false
       else:
         conf_matrix = test(E+, E-, B, hyp)
         outcome = compute_outcome(conf_matrix)
         score = compute_score(conf_matrix)
         if outcome == ('all', 'none'):
            best_hyp = hyp
            end_popper = true
         else:
            if best_score < score: 
               best_hyp = hyp
            cc += learn_constraints(hyp, outcome)

   return best_hyp
\end{lstlisting}
\end{algorithm}









